\section{Parte 6: detrás de la serie hay clientes}
\label{sec:parte6}

El análisis RFM usa \texttt{limpiar(df, requiere\_cliente=True)}, es decir,
las reglas 1 a 6 de la bitácora: las mismas de la serie más la que quita las
filas sin \texttt{CustomerID}. Todo lo genera \texttt{parte6.py}. Para cada
cliente se calcula:
\begin{itemize}
  \item \textbf{R}: días entre su última factura de venta (no cancelación) y
    el 1 de octubre de 2011;
  \item \textbf{F}: número de facturas de venta distintas;
  \item \textbf{M}: venta neta total, con las cancelaciones restando.
\end{itemize}

\subsection{Inciso 6.1: R, F y M antes del 1 de octubre}

Con las transacciones anteriores al 1 de octubre de 2011 hay \textbf{3\,658
clientes}. La tabla~\ref{tab:p6_resumen} resume sus tres variables, que son
muy asimétricas: la mitad de los clientes compró dos veces o menos, mientras
que el que más compró tiene 139 facturas, y el monto va de $-$\pounds4\,288 a
\pounds201\,772, con mediana de \pounds562.

\begin{table}[H]
  \centering
  \small
  \tabla{p6_rfm_resumen}
  \caption{Distribución de R, F y M para los 3\,658 clientes con transacciones
  antes del 1 de octubre de 2011. P20, P40, P60 y P80 son los cortes de los
  quintiles.}
  \label{tab:p6_resumen}
\end{table}

\paragraph{Clientes con $M\leq0$: 54.} De ellos, 51 tienen $M<0$ y 3 tienen
$M=0$. Hay dos tipos:
\begin{itemize}
  \item \textbf{43 no tienen ninguna factura de venta} en el periodo, solo
    cancelaciones: devolvieron mercancía que habían comprado antes del 1 de
    diciembre de 2010, fuera de los datos. No tienen una ``última compra''
    observada, así que se les asigna $R=305$, un día más que a quien compró
    el primer día de los datos (su última compra fue antes).
  \item \textbf{11 compraron y luego cancelaron} lo mismo o más, como el
    cliente 12346 del inciso 1.4 (74\,215 frascos vendidos y cancelados el
    mismo día, $F=1$ y $M=0$).
\end{itemize}
\emph{Se quedan en el análisis}, como pide el inciso 6.2, y la decisión es
defendible: son clientes reales que tuvieron actividad, y su $M$ (y, para los
43, su $R$) los manda naturalmente al puntaje más bajo. Quitarlos cambiaría
los cortes de los quintiles de todos los demás. Lo que no se hace es tratarlos
como clientes valiosos: en ningún análisis de valor pueden aparecer arriba.

\paragraph{¿Por qué no usar nada posterior al 1 de octubre?} Porque lo que se
quiere evaluar es si R, F y M \emph{anticipan} quién compra entre el 1 de
octubre y el 8 de diciembre. Si se usaran transacciones posteriores, la
respuesta se colaría en las variables: un cliente que compró en noviembre
tendría una recencia pequeña \emph{porque} compró en noviembre, y una
frecuencia y un monto mayores por la misma razón. El AUC del inciso 6.3 saldría
excelente y no mediría nada. Es la misma fuga del inciso 5.4: una variable
que contiene el futuro hace que cualquier modelo se vea bien.

\subsection{Inciso 6.2: puntajes por quintiles}

A cada cliente se le asigna un puntaje de 1 a 5 en cada dimensión según su
quintil; en recencia, 5 es el más reciente. Se siguió exactamente el
procedimiento del enunciado: se ordenan todos los clientes por
\texttt{CustomerID}, se barajan con
\texttt{numpy.random.default\_rng(2026).permutation} y se aplica
\texttt{pd.qcut(x.rank(method='first'), 5)}, con $x=F$, $x=M$ y $x=-R$.

\paragraph{El problema con F.} F es un conteo pequeño con muchísimos empates:
el 39.6\% de los clientes tiene $F=1$, el 18.9\% tiene $F=2$, el 12.1\% tiene
$F=3$ y el 1.2\% tiene $F=0$. El primer quintil (20\%) y el segundo (40\%)
caen dentro del bloque de $F=1$, así que los cortes de los quintiles se
repiten: \texttt{pd.qcut(F, 5)} directamente falla con el error ``Bin edges
must be unique'', porque los cortes son $0, 1, 1, 3, 5, 139$. No hay forma de
partir a los clientes en cinco grupos del mismo tamaño usando solo su F.

\begin{table}[H]
  \centering
  \small
  \tabla{p6_puntaje_f}
  \caption{Número de clientes por valor de $F$ y puntaje asignado, para
  $F\leq6$. Los clientes con el mismo $F$ que aparecen en varias columnas
  recibieron su puntaje por sorteo.}
  \label{tab:p6_puntaje_f}
\end{table}

\paragraph{Lo que implica el desempate.} Con \texttt{rank(method='first')},
los clientes empatados se ordenan según su posición en el barajado, así que el
puntaje de los clientes que caen en una frontera se decide por sorteo
(tabla~\ref{tab:p6_puntaje_f}). Los 1\,449 clientes con $F=1$ quedaron
repartidos entre el puntaje 1 (689), el 2 (731) y el 3 (29); los de $F=3$,
entre 3 y 4; los de $F=5$, entre 4 y 5. Para dos clientes con la misma F que
caen en puntajes distintos, \emph{la diferencia de puntaje no significa nada}:
se comportaron igual y uno sacó mejor número en la rifa. Con otra semilla los
puntajes cambian, y cualquier análisis que compare los puntajes 1 y 2 de F
está comparando dos muestras al azar del mismo grupo. Se nota en el inciso
6.3: los puntajes 1 y 2 de F tienen exactamente la misma tasa de compra (31.8\%
y 31.7\%). Fijar la semilla hace que todos obtengan los mismos puntajes, pero
no les da significado.

\subsection{Inciso 6.3: ¿quién compró del 1 de octubre al 8 de diciembre?}

Sea $\text{compra}_i=1$ si el cliente $i$ tiene alguna factura de venta entre
el 1 de octubre y el 8 de diciembre de 2011. Lo hizo el 50.4\% de los 3\,658
clientes. La tabla~\ref{tab:p6_tasas} y la figura~\ref{fig:p6_tasas} muestran
la tasa de compra por puntaje en cada dimensión.

\begin{table}[H]
  \centering
  \small
  \tabla{p6_tasas}
  \hspace{2em}
  \tabla{p6_auc}
  \caption{Izquierda: fracción de clientes que compró entre el 1 de octubre y
  el 8 de diciembre, por puntaje. Derecha: AUC de cada variable cruda y de la
  suma de los tres puntajes como clasificadores de $\text{compra}_i$.}
  \label{tab:p6_tasas}
\end{table}

\begin{figure}[H]
  \centering
  \includegraphics[width=0.75\textwidth]{p6_tasas.png}
  \caption{Tasa de compra en la ventana de prueba por puntaje de R, F y M.}
  \label{fig:p6_tasas}
\end{figure}

Las tres dimensiones ordenan bien: la tasa de compra crece con el puntaje en
las tres. Pero no separan igual.
\begin{itemize}
  \item \textbf{F separa mejor}: del 31.7\% en los puntajes bajos al 85.7\% en
    el puntaje 5, el rango más amplio, y el mayor AUC de las tres variables
    crudas, 0.721. Quien ya compró muchas veces vuelve a comprar.
  \item \textbf{M} queda muy cerca (de 28.8\% a 80.2\%, AUC 0.710). Es
    natural: un cliente que compra muchas veces acumula más monto, así que M
    y F miden en buena parte lo mismo.
  \item \textbf{R} es la que menos separa (AUC 0.674). Distingue bien entre
    los puntajes bajos (27.5\% en el 1, 55.2\% en el 3), pero se satura
    arriba: los puntajes 4 y 5 tienen la misma tasa (65.9\% y 66.1\%). Haber
    comprado hace dos semanas o hace seis no hace diferencia en una ventana
    de diez semanas que además es la temporada alta, cuando incluso los
    clientes menos recientes vuelven.
  \item \textbf{La suma de los tres puntajes} tiene el mejor AUC, 0.735,
    pero apenas por encima de F sola: combinar las tres dimensiones agrega
    poco, porque están muy correlacionadas.
\end{itemize}
Un AUC de 0.72 significa que, si se toma al azar un cliente que compró y uno
que no, F le da el valor más alto al que compró el 72\% de las veces: útil,
pero lejos de una predicción segura.

\subsection{Inciso 6.4: los campeones}

Los clientes con puntaje 4 o 5 en las tres dimensiones son \textbf{792, el
21.7\% de los clientes}. Entre el 1 de octubre y el 8 de diciembre
\textbf{generaron el 53.5\% de la venta neta} de todos los clientes con
\texttt{CustomerID}, incluidos los nuevos: \pounds1.28 millones de
\pounds2.40 millones (tabla~\ref{tab:p6_grupos}). El 81.2\% de los campeones
compró en esa ventana. Es la concentración típica de un negocio mayorista: uno
de cada cinco clientes trae más de la mitad del dinero. Para una tienda, saber
quiénes son vale más que cualquier pronóstico de la serie agregada.

\begin{table}[H]
  \centering
  \small
  \tabla{p6_grupos}
  \caption{Reparto de la venta neta del 1 de octubre al 8 de diciembre de 2011
  entre los clientes con \texttt{CustomerID}. ``Clientes con movimiento'' es
  cuántos clientes de cada grupo tuvieron alguna transacción en la ventana.
  Un cliente es nuevo si no tiene ninguna factura de venta antes del 1 de
  octubre.}
  \label{tab:p6_grupos}
\end{table}

\subsection{Inciso 6.5: los clientes nuevos}

De los 2\,596 clientes con movimiento entre el 1 de octubre y el 8 de
diciembre, \textbf{723 son nuevos}: no tienen ninguna factura de venta antes
del 1 de octubre. Generaron \textbf{el 15.7\% de la venta neta} de la ventana,
\pounds377\,860 (tabla~\ref{tab:p6_grupos}).

\paragraph{¿Puede un RFM pronosticar esa venta?} No, ni el RFM ni ningún
modelo construido cliente por cliente. Un modelo así solo puede pronosticar a
los clientes que tiene en su tabla, y un cliente nuevo no tiene fila: no hay R,
ni F, ni M que calcular, porque todavía no ha hecho nada. Lo más que puede
hacer un modelo por cliente es predecir que los clientes conocidos dejen de
comprar; nunca que aparezcan clientes desconocidos. Ese 15.7\% es un error
estructural, que no se corrige con un mejor modelo. Y es un mínimo: la venta
sin \texttt{CustomerID}, el 15.4\% de toda la venta según el inciso 1.2(f),
tampoco la ve ningún modelo por cliente.

\paragraph{Conexión con la Parte 5 y con el Proyecto 1.} En el vocabulario de
FPP2 (§10.3--10.4), la Parte 5 es un pronóstico \emph{top-down}: se
pronostica directamente la venta total. Un modelo que sume pronósticos cliente
por cliente sería \emph{bottom-up}. Solo el \emph{top-down} puede ver a los
clientes nuevos. No los ve como individuos, sino como parte del nivel de la
serie: la llegada de cientos de clientes nuevos en el cuarto trimestre es una
de las razones por las que la venta se duplicó desde septiembre, y es
justamente ese cambio de nivel el que S2 siguió y le dio la victoria en el
inciso 5.3. El \emph{bottom-up}, en cambio, empieza con una desventaja: aunque
pronosticara perfecto a cada cliente conocido, se quedaría corto en al menos
un 15.7\% de la venta con cliente. Para la comparación que pide el Proyecto 1,
esto significa que un pronóstico cliente por cliente necesita un componente
aparte para la venta de clientes nuevos (que tendría que ser de algún modo
\emph{top-down}), o va a perder contra la serie agregada en cualquier periodo
en que la tienda gane clientes, como en la temporada navideña.
